There is a sense in the previous Chapters that we want to study some distributions when our matrices are large enough - be that for feasibility of computations or approximation of infinite dimension operators. We also have mentioned that our main interest here is the Invariant Ensembles, or specially in Orthogonal Polynomial ones. Because of this, studying large matrices is equivalent, in some sense, to studying the regime of large $n$ for orthogonal polynomials. For some of those polynomials, a integral representation is available, allowing us to use asymptotic methods on integrals to solve for the large $n$ behavior. This is the main discussion of this chapter.

Techniques and structuring here are based on the book by \citeonline{NISTMathFunc} and \citeonline{Trapezoidal-Maclaurin} for quadrature and on the lecture notes by \citeonline{guilhermeasympt} for asymptotics on integrals. 

\section{Quadrature Technique}

Before doing asymptotics on integrals we discuss a technique that will become important when approximating large summations: quadratures, i.e., approximating integrals by summations and vice-versa. One can see that asymptotics on integrals can then also perform a role on discrete summations. We first start briefly introducing a useful tool in the development of the formula we wish to use: the Bernoulli numbers $B_n$, polynomials $B_n(x)$ and periodic polynomials $\tilde{B}_n(x)$. The Bernoulli numbers $B_n$ are a sequence of rational numbers that frequently appears in the Taylor series expansion of trigonometric-type functions. These number can also be seen as special values obtained by the Bernoulli polynomials $B_n(x)$. This is because these polynomials can be written explicitly as
\begin{equation}
	B_k(x) = \sum_{n=0}^{k} {k \choose n} B_n x^{k-n},
\end{equation}
then, we must have that $B_k(0) = B_k$. More interestingly for us is the fact that these polynomials satisfy the differential equation 
\begin{equation}
	B'_n(x) = n B_{n-1}(x), \quad n = 1,2,\cdots
	\label{Equation: Differential Bernoulli}
\end{equation}
Finally, the periodic Bernoulli polynomial is simply the Bernoulli polynomials evaluated at the fractional part of the argument, i.e., $\tilde{B}_n(x) = B_n(x - \floor{x})$. Now, let $B_n$ and $B_n(x)$ denote the $n$th Bernoulli number and polynomial, respectively, and $\tilde{B}_n(x)$ the $n$th Bernoulli periodic function $B_n(x - \floor{x})$. Assume that $a, m$ and $n$ are integers such that $n > a$, $m > 0$, and $f^{(2m)}(x)$ is absolutely integrable over $[a,n]$. Then, 

\begin{teorema} (Euler-Maclaurin Formula,  \cite{NISTMathFunc})
	Following the context given above, 
	\begin{multline}
		\sum_{j=a}^{n} f(j) = \int_{a}^{n} f(x) \dd x + \frac{1}{2} \left(f(a) + f(n)\right) + \sum_{s=1}^{m-1} \frac{B_{2s}}{(2s)!} \left( f^{(2s-1)}(n) - f^{(2s-1)}(a)\right) \\ + \int_{a}^{n} \frac{B_{2m} - \tilde{B}_{2m}(x)}{(2m)!} f^{(2m)}(x) \dd x.
	\end{multline}
\end{teorema}

This identity is called the \textit{Euler-Maclaurin} formula. It can be extremely useful when estimating the asymptotic of the summation on the function $f$ when its integral can be more easily evaluate. In essence, this is a quadrature technique; a numerical integration: the approximation of an integral $\int f(x) \dd x$ of some function $f$ by a discrete summation $\sum w_i f(x_i)$ over a set of points $x_i$ with some weight function $w_i$. For example, in the classical Riemann sums, one approximates the integral $\int f(x) \dd x$ over some interval $[a,b]$ with the limit of the summation $\sum f(x_i^*) \Delta x_i$, where $x_i^* \in [x_{i-1}, x_i]$ and $ \Delta x_i = x_i - x_{i-1}$. Different methods can be generated by choosing $x_i^*$ differently.

The trapezoidal rule, for example, is a Riemann sum where $x_i^*$ is chosen to be the mean point $x_i^* = (x_{i-1} + x_i)/2$. In this way, we would estimate the area using trapezoids, as the name suggests. Then, one could write
\begin{align}
	\int_{a}^{n} f(x) \dd x & \approx \sum_{j = 0}^{n} f(x_j^*) = \sum_{j = 0}^{n} \frac{f(x_{i} + x_{i-1})}{2}, \\
	& = \sum_{j = 1}^{n-1} f(x_j^*) + \frac{1}{2} \left( f(n) + f(0) \right) = \sum_{j = 0}^{n-1} f(x_j^*) + \frac{1}{2} \left( f(n) - f(n) \right) .
\end{align}
Of course, this has some error term. Suppose that for some $(x_i, x_{i-1})$ we have the same $k$-derivatives at both extreme points for every $k\geq 1$, i.e., the function is constant. Then, there is no error in the approximation. If however, the two extreme points differ only on the $k$th derivative, there is some curvature defined by this difference where the error comes from. Therefore, there is indication of some error function given in respect to higher order derivatives.

\begin{proof}
Following the computations in \citeonline{Trapezoidal-Maclaurin}, take a function $f \in C^p(\R)$ for $p \geq 1$. To derive the Euler-Maclaurin for the interval $[a,b]$, let $h = (b - a)/n$ and $c = a + j h$, then we express
\begin{equation}
	\int_{a}^{b} f(x) \dd x = \sum_{j=0}^{n-1} \int_{a+j h}^{a+(j+1)h} f(x) \dd x
\end{equation}
where we use \eqref{Equation: Differential Bernoulli} and integration by parts to rewrite each integral as
\begin{align}
	\int_{c}^{c+h} f(x) \dd x & = h \int_{0}^{1} B_0(1-x) f(c+xh) \dd x \\
	& = - h \frac{B_1(1-x)}{1!} f(c + xh) |_0^1 + h^2 \int_{0}^{1} \frac{B_1(1-x)}{1!} f'(c+xh) \dd x, \\
	& \vdots \nonumber \\
	& = - \sum_{r = 0}^{p-1} \frac{h^{r+1} B_{r+1}(1-x)}{(r+1)!} f^{(r)}(c+xh) |_0^1 + h^{p+1} \int_0^1 \frac{h^{r+1} B_{p}(1-x)}{(p)!} f^{(p)}(c+xh) \dd x, \\
	& 
	\begin{multlined}
		= h \frac{f(c) + f(c+h)}{2} - \sum_{r = 0}^{p-1} \frac{h^{r+1} B_{r+1}(1-x)}{(r+1)!} \left[ f^{(r)}(c+h) - f^{(r)}(c) \right] \\
		+ h^{p+1} \int_0^1 \frac{B_{p}(1-x)}{(p)!} f^{(p)}(c+xh) \dd x. 
	\end{multlined}
\end{align}
where one can start to notice, for instance, the trapezoidal rule creeping out on the first term. Notice that it was necessary to use that $ - B_1(0) = B_1(1) = 1/2$ and $B_n(0) = B_n(1) = B_n$ for $n\neq1$. Returning to the summation, we rearrange terms as
\begin{align}
	h & \left[ \frac{f(a)}{2} + f(a+h) + \cdots + f(b-h) + \frac{f(b)}{2} \right] = \frac{1}{2} (f(b) - f(a)) + \sum_{j=0}^{n-1} f(a + j h), \\
	& = \int_{a}^{b} f(x) \dd x + \sum_{r = 0}^{p-1} \frac{h^{r+1} B_{r+1}(1-x)}{(r+1)!} \left[ f^{(r)}(b) - f^{(r)}(a) \right] - h^{p+1} \int_0^1 \frac{B_{p}(1-x)}{(p)!} f^{(p)}(c+xh) \dd x. 
\end{align}
\end{proof}

\section{Asymptotic on Integrals}

The techniques we will expand on here refers to a specific type of complex function $\phi: \C \rightarrow \R$ in the regime where $\lambda \rightarrow \infty$. Those functions are of the form 
\begin{equation}
	\phi(\lambda) := \int_{\Gamma} g(s) \ee^{-\lambda f(s)} \dd s,
	\label{Equation: Integral form phi}
\end{equation}
where $\Gamma$ is a given contour in the plane and $f, g$ are suitable functions. Suitable here means that we impose $g$ and $f$ to be continuous (or analytic); $g$ is integrable; $f$ have global minima with zero derivative and strictly positive second derivative at this point of minima. There is also a more technical regularity condition on $f$ to ensure the approximations hold to still be made clear.

\begin{exemplo}
	(Airy Equation) Function such as in \eqref{Equation: Integral form phi} are not uncommon. Consider the function that satisfies the equation $\psi''(x) = x \psi(x)$. Taking the Fourrier transform we have 
	\begin{equation}
		-\omega^2 \Phi(x) = \ii \Phi'(x)
	\end{equation}
	for which we solve with $\Phi(x) = \ee^{- \ii \omega^3 / 3}$ and then take the inverse Fourrier Transform to get 
	\begin{equation}
		\phi(x) = \frac{1}{2\pi} \int_{-\infty}^{\infty} \ee^{\ii/3 \omega^3 + \ii x \omega} \dd \omega, \quad x \in \R.
	\end{equation}
 	The imaginary part of the argument is an odd function and therefore zero when integrated. The real part is expressed 
	\begin{equation}
		\phi(x) = \frac{1}{2\pi} \int_{-\infty}^{\infty} \cos{\left( \frac{\omega^3}{3} + x \omega \right)} \dd \omega.
	\end{equation}
	This is the Airy function. For this integral, the techniques applied in this chapter are of great use to understand its behavior as $\omega \rightarrow \infty$.
\end{exemplo}

Moreover, one can somehow somehow motivate the conditions we imposed on $\phi$. Consider $\phi: \R \rightarrow \R$, $f$ with a single global minima and $\Gamma \equiv \R$ - for simplicity; the argument's idea should hold for other choices given the appropriate considerations. If $\phi$ is as in \eqref{Equation: Integral form phi}, and satisfying the conditions given for such function, then take $x^*$ the single minimal critical point for $f$. We first note that if that is true, then, for any $|\epsilon| > 0$, the following limit should hold:
\begin{equation}
	\lim_{\lambda \rightarrow \infty} \frac{\ee^{-\lambda(f(x^*))}}{\ee^{-\lambda(f(x^* - \epsilon))}} = \lim_{\lambda \rightarrow \infty} \ee^{-\lambda(f(x^*) - f(x^* - \epsilon))} = \lim_{\lambda \rightarrow \infty} \ee^{\lambda |c|} \rightarrow + \infty. 
\end{equation}
That is, the main contribution for the integral comes from a single point. The main idea then is to separate the integral in two regions: A first interval $(x^* - c/\sqrt{\lambda}, x^* + c/\sqrt{\lambda}), c > 0$, including this critical point and another interval that in theory should contribute only in error for the integral asymptotic. We start by approximating $f$ around its global maximum by Taylor's theorem 
	\begin{equation}
		f(x) = f(x^*) - \frac{1}{2} |f''(x^*)|(x-x^*)^2 + \epsilon(x-x^*)
	\end{equation}
	and then perform a change of variables $\frac{x}{\sqrt{\lambda}} = s - x^*$, so that 
	\begin{equation}
		\int_{x^*-\frac{c}{\sqrt{\lambda}}}^{x^*+\frac{c}{\sqrt{\lambda}}} g(s) \ee^{-\lambda f(s)} \dd(s) = \frac{\ee^{-\lambda f(x^*)}}{\sqrt{\lambda}} \int_{-c}^c g\left(x^* + \frac{x}{\sqrt{\lambda}}\right) \ee^{\frac{1}{2} |f''(x^*)|x^2 - \epsilon \left(\frac{x^3}{\sqrt{\lambda}}\right)} \dd x.
	\end{equation}
	Now, using the technical assumption of regularity on the function $f$, one can argue that the error term $\epsilon$ goes to zero sufficiently fast, leaving a Gaussian Integral. Moreover, one needs to argue that the integral outside this small neighborhood of $x^*$ is exponentially smaller then $\ee^{-\lambda f(x^*)}$. This should follow by an argument where the asymptotic for the integral 
	\begin{equation}
			\int_{x^*+ \frac{c}{\sqrt{\lambda}}}^{\infty} |g(s)| \ee^{-\lambda f(s)} \dd(s) \leq \ee^{-\lambda f(x^*+ c/\sqrt{\lambda})} \normL{1}{g}{}
	\end{equation}
	is exponentially smaller than the previous region. 
	\begin{figure}[h!]
		\centering
		\begin{tikzpicture}
			\begin{axis}[
				width=\textwidth, height=8cm,
				xmin=-3,ymin=0,ymax=1.1,xmax=3
				]
				  \path[name path=axis] (axis cs:-3,0) -- (axis cs:3,0);
				\addplot[name path=f1, domain=-3:3,smooth] {e^(-1*(cosh(x)-1))} ;
				\addplot[name path=f2, domain=-3:3,smooth] {e^(-2*(cosh(x)-1))} ;
				\addplot[name path=f3, domain=-3:3,smooth] {e^(-10*(cosh(x)-1)};
				\addplot[name path=f4, domain=-3:3,smooth] {e^(-70*(cosh(x)-1))};
				\addplot[name path=f5, domain=-3:3] {e^(-1000*(x^2))};
				\addplot [thick,color=black,fill=black, fill opacity=0.05] fill between[of=f1 and axis,soft clip={domain=-3:3}];
				\addplot [thick,color=black,fill=black, fill opacity=0.05] fill between[of=f2 and axis,soft clip={domain=-3:3}];
				\addplot [thick,color=black,fill=black, fill opacity=0.05] fill between[of=f3 and axis,soft clip={domain=-3:3}];
				\addplot [thick,color=black,fill=black, fill opacity=0.05] fill between[of=f4 and axis,soft clip={domain=-3:3}];
				\tikzfillbetween[of=f5 and axis]{pattern=dots};% <- changed
			\end{axis}
		\end{tikzpicture}
		\caption{In the Plot it is represented the plot for $\ee^{- \lambda (\cosh(x)-1)}$ and an increasing $\lambda  = 1, 2, 10, 70$ - darker represent larger values of the parameter. The dotted patter represents $\ee^{-\lambda x^2}$ for large lambda.}
	\end{figure}
		
	With this, we only note a few details of interest. Firstly, more precise error estimates can be obtained if one uses higher order terms of the Taylor Series expansion of $f$. For maximizer $x^*$ which have null second derivative (a double critical point), the Laplace Approximation can be used expanding to the third order terms of the Taylor Series $f(x) = f(x^*) = \frac{1}{6}f'''(x^*)(x-x^*)^3 + \epsilon(x-x^*)$ and making the change of variables $t/\sqrt[3]{n} = x - x^*$ instead. This argument and result could also be used on curves on the plane. This is motivated from the fact that every (smooth) curve is locally a straight line. This will lead us to the method of \textit{Steepest Descent}. These next subsections are in the spirit of proving a bit more formally the expansion given by the \textit{Laplace Method} and introducing the mentioned \textit{Steepest Descent}.

\subsection{Watson's Lemma} 

We start this topic on the real line by looking at the functions $F: \R \rightarrow \R$ of the form
\begin{equation}
	F(\lambda) = \int_{0}^{T} f(t) \ee^{-\lambda t} \dd t,
	\label{equation: F}
\end{equation}
which one may think of as the Laplace transform of a function $f$ supported on $[0,T]$. Note that $F$ fits well with the models we described previously, and so one should expect the main contribution to come from the minima of the exponent function. Let us formalize the conditions for such expectation and the consequent result.

\begin{teorema}	
	(Watson's Lemma) Take $f \in L^1(0,T)$ and suppose that for some $\delta > 0$ it is of the form
	\begin{equation}
		f(t) = t^\alpha f_0(t), \quad 0 < t \leq \delta,
	\end{equation}
	where $\alpha \in (-1, \infty)$ and $f_0$ is continuously differentiable on $[0, \delta]$ with $f_0(0) \neq 0$. Then the function $F$ in \eqref{equation: F} satisfies
	\begin{equation}
		F(\lambda) = \frac{f_0(0) \Gamma(\alpha + 1)}{\lambda^{\alpha + 1}} (1 + \Boh(\lambda^{-1}))
	\end{equation}
	when $\lambda \rightarrow \infty$. Also, $\Gamma$ is the usual Gamma function.
	\label{Theorema: Watson}
\end{teorema}

We must first mention that a weak estimation of the integral could be given by
\begin{equation}
	\int_{0}^{T} f(t) \ee^{-\lambda t} \dd t \leq \normL{1}{f}{(0, T)}
\end{equation}
where we only use the fact that $f \in L^1(0,T)$. However, we have more information on the behavior of the function in the section $0 < t < \delta$. Because of that fact, it is expected that in this set we could refine the approximation.

\begin{proof}
	 We start by decomposing the function in two parts
	\begin{equation}
		F(\lambda) = \int_{0}^{\delta} f(t) \ee^{-\lambda t} \dd t + \int_{\delta}^{T} f(t) \ee^{-\lambda t} \dd t \revdeff (1) + (2).
	\end{equation}
	For $(2)$, we use the same weak estimate as commented before,
	\begin{equation}
		\int_{\delta}^{T} f(t) \ee^{-\lambda t} \dd t \leq  \ee^{-\lambda \delta} \int_{\delta}^{T} |f(t)| \dd t = \Boh(\ee^{-\lambda \delta}  \normL{1}{f}{(\delta, T)}).
	\end{equation}
	We keep it this way as we have no further information on the $f$ in this section of the real line. This leave us with 
	\begin{equation}
		F(\lambda) = \int_{0}^{\delta} f(t) \ee^{-\lambda t} \dd t + \Boh(\ee^{-\lambda \delta}  \normL{1}{f}{(\delta, T)}).		
	\end{equation}
	Now, we remind of the condition of continuous differentiability and the hypothesis on $f(t)$ to get $f(t) = t^\alpha (f_0(0) + r(t))$ for some differentiable error function $r(t)$ satisfying $|r(t)| \leq \normL{\infty}{f_0'}{(0, \delta)} t$ with $\normL{\infty}{f'_0}{(0, \delta)}  = \inf{\{ c \geq 0 \ | \ \mu|_{(0, \delta)}(\{|f'_0| > c\}) = 0 \}}$ as usual. Then, write $(1)$ as
	\begin{align}
		(1) = \int_{0}^{\delta} t^\alpha(f_0(0) + r(t)) \ee^{-\lambda t} \dd t & = f_0(0) \int_{0}^{\delta} t^\alpha \ee^{-\lambda t} \dd t + \int_{0}^{\delta} t^\alpha r(t) \ee^{-\lambda t} \dd t.
	\end{align}
	For these two integrals we simplify the problem by rewriting 
	\begin{align}
	\int_{0}^{\delta} t^\alpha \ee^{-\lambda t} & = \int_{0}^{\infty} t^\alpha \ee^{-\lambda t} \dd t - \int_{\delta}^{\infty} t^\alpha \ee^{-\lambda t} \dd t  = \frac{\Gamma(\alpha + 1)}{\lambda^{\alpha +1}} - \int_{\delta}^{\infty} t^\alpha \ee^{-\lambda t} \dd t, \\
	\left|\int_{0}^{\delta} t^\alpha r(t) \ee^{-\lambda t} \dd t \right| & \leq \normL{\infty}{f_0'}{(0, \delta)} \int_{0}^{\delta} t^{\alpha + 1} \ee^{-\lambda t} \dd t, \\ &=  \normL{\infty}{f_0'}{(0, \delta)} \left( \int_{0}^{\infty} t^{\alpha+1} \ee^{-\lambda t} \dd t - \int_{\delta}^{\infty} t^{\alpha+1} \ee^{-\lambda t} \dd t \right), \\
	& =  \normL{\infty}{f_0'}{(0, \delta)} \left( \frac{\Gamma(\alpha + 2)}{\lambda^{\alpha +2}} - \int_{\delta}^{\infty} t^{\alpha+1} \ee^{-\lambda t} \dd t \right).
	\end{align}
	We are now only missing one ingredient to get the result, the order of
	\begin{equation}
		G(\lambda, \beta) \deff \int_{\delta}^{\infty} t^{\beta} \ee^{-\lambda t} \dd t, \ \ \ \lambda > 1, \quad \beta > -1. 
	\end{equation}
	Using Cauchy-Schwartz we estimate
		\begin{align}
			|G(\lambda, \beta)|^2 & \leq \left( \int_{\delta}^{\infty} \ee^{-\lambda t} \dd t \right) \left(\int_{\delta}^{\infty} t^{2\beta} \ee^{-\lambda t} \dd t \right), \\
			& = \left(\frac{e^{-\lambda \delta}}{\lambda}\right) \left( \int_{\delta}^{\infty} t^{2\beta} \ee^{-\lambda t}\dd t \right),	\\
			& \stackrel{(u=\lambda t)}{=} \left(\frac{e^{-\lambda \delta}}{\lambda}\right) \left( \int_{\delta}^{\infty} \ee^{-u} u^{2\beta} \frac{1}{\lambda^{2\beta + 1}} \dd u \right), \\
			& = \frac{\ee^{-\lambda \delta}}{\lambda^{2\beta + 2}} M_\beta
		\end{align}
	where $M_\beta$ is some constant independent of $\lambda$. That leave us with
	\begin{equation}
		G(\lambda, \beta) = \Boh\left(\frac{\ee^{-\lambda\delta/2}}{\lambda^{\beta+1}}\right) = \Boh\left(\ee^{-\epsilon\lambda}\right)
	\end{equation}
	where $\epsilon > 0$ depends on $\beta$ and $\delta$. Finally we rewrite
	\begin{equation}
		\int_{0}^{\delta} t^\alpha \ee^{-\lambda t}  = \frac{\Gamma(\alpha + 1)}{\lambda^{\alpha +1}} - G(\lambda, \alpha) =  \frac{\Gamma(\alpha + 1)}{\lambda^{\alpha+1}} + \Boh\left(\ee^{-\epsilon\lambda}\right),
	\end{equation}
	\begin{align}
		\left|\int_{0}^{\delta} t^\alpha r(t) \ee^{-\lambda t} \dd t \right| & \leq \normL{\infty}{f'}{(0, \delta)} \left( \frac{\Gamma(\alpha + 2)}{\lambda^{\alpha +2}} - G(\lambda, \alpha + 1) \right) \\
		&  =  \normL{\infty}{f'}{(0, \delta)} \left( \frac{\Gamma(\alpha + 2)}{\lambda^{\alpha +2}} + \Boh\left(\ee^{-\epsilon\lambda}\right) \right) = \Boh\left(\lambda^{-\alpha-2}\right),
	\end{align}
	and
	\begin{equation}
		F(\lambda) = \Boh\left(\lambda^{-\alpha-2}\right) + f_0(0) \left( \frac{\Gamma(\alpha + 1)}{\lambda^{\alpha+1}} + \Boh\left(\ee^{-\epsilon\lambda}\right) \right) + \normL{1}{f}{(\delta, T)} \ee^{-\delta \lambda}.
	\end{equation}
	By the dominance of the exponential error, we get the result
	\begin{equation}
		F(\lambda) = \frac{f_0(0) \Gamma(\alpha + 1)}{\lambda^{\alpha + 1}} (1 + \Boh(\lambda^{-1})).
	\end{equation}
\end{proof}

In the estimation above, the Gamma function and the factor of $\lambda^{-1}$ is the result of integrating the linear exponent of $\ee^{-\lambda t}$. If we, however, change this exponent to be of the general type $-\lambda t^d$ we still have a similar result on the behavior of the function.

\begin{proposicao}		\label{prop: cor watson}
	For $T, S > 0$, define the function
	\begin{equation}
		G(\lambda) \deff \int_{-S}^{T} g(t) \ee^{-\lambda t^2} \dd t
	\end{equation}
	and assume $g(x)$ absolutely integrable in $[-S, T]$, of class $C^1$ in a fixed neighborhood of the origin. Then the function satisfies
	\begin{equation}
		G(\lambda) = \frac{g(0) \Gamma(1/2)}{\lambda^{1/2}} + \Boh \left(\frac{1}{\lambda^{3/2}} \right)
	\end{equation}
	when $\lambda \rightarrow \infty$ and where $\Gamma$ is the Gamma function.
\end{proposicao}

\begin{proof}
	The proof is a reduction of the problem to Watson's Lemma. Assume, $T>S$ and write
	\begin{equation}
		G(\lambda) = \int_{-S}^{S} g(t) \ee^{-\lambda t^2} \dd t + \int_{S}^{T} g(t) \ee^{-\lambda t^2} \dd t.
	\end{equation}
	The second integral we know to be of order
	\begin{equation}
		\int_{S}^{T} |g(t)| \ee^{-\lambda t^2} \dd t \leq ||g||_{\infty} \ee^{-\lambda S^2},
	\end{equation}
	exponential on $\lambda$, so we drop it. For the first integral we rewrite
		\begin{align}
			\int_{-S}^{S} g(t) \ee^{-\lambda t^2} \dd t & = \int_{0}^{S} (g(t) + g(-t)) \ee^{-\lambda t^2}\dd t, \\
			& \stackrel{(u= t^2)}{=} \frac{1}{2} \int_{0}^{S^2} \left(\frac{g\left(\sqrt{u}\right) + g\left(-\sqrt{u}\right)}{\sqrt{u}}\right) \ee^{-\lambda u} \dd u
		\end{align}
	where now we use Watson's Lemma to finish the proof.
\end{proof}

\begin{exemplo}
	Take the function called Gamma given by the expression
	\begin{equation}
		\Gamma(x) = \int_{0}^{\infty} \ee^{-t} t^{x-1} \dd t.
		\label{eq: Gamma}
	\end{equation}
	We are interested in what happens to the function when we take $x \rightarrow +\infty$. To see what happens, let us fit it into \autoref{Theorema: Watson}. Rewrite \eqref{eq: Gamma} as
	\begin{equation}
		\int_{0}^{\infty} \ee^{-t+x\log{(t)}-\log{(t)}} \dd t.
	\end{equation} 
	This lead us to identify $\phi$ as $-\log{(t)}$, which would give us $\phi'(t) < 0$, that is, we would have its minimum at $t=+\infty$. However $\ee^{\log{(t)}} = t^1$ is not integrable near the origin. Something need to change about the identification. Let us introduce a variable to rescale the integral: $xs = t$, s.t.,
	\begin{equation}
		\Gamma(x) = x^x \int_{0}^{\infty} \ee^{-\lambda(s - \log{(s)})} \ee^{-s} \dd s,
	\end{equation}
	where we have used $\lambda = x - 1$. Thus we identify $\phi(s) = s - \log(s)$ and $g(s) = \ee^{-s}$. In fact, now $\phi$ has a global minimum at $s = 1$ and $\phi''(1) = 1 \neq 0$. Hence
	\begin{align}
		\Gamma(x) & = \frac{\sqrt{2\pi} x^x \ee^{-x}}{(x-1)^{(-1/2)}} (1 + \Boh(x^{-1})), \\
		& = \sqrt{2\pi} x^{x + 1/2} \ee^{-x} (1 + \Boh(x^{-1})),
	\end{align}
	which is called the Stirling formula.
\end{exemplo}

\subsection{Laplace Method}

A seemly extension of the type of function we were dealing with are functions of the type
\begin{equation}
	\Phi(\lambda) \deff \int_{a}^{b} g(t) \ee^{-\lambda \phi(t)} \dd t
	\label{equation: Laplace}
\end{equation}
for $a,b \in \R \cup \{-\infty, \infty\}$ and $a<b$ and a well behaved function $g$. For now, $\phi$ is a real-valued function. To begin the study let us make some assumptions. The first one is that $a$ is finite, that means that we set $a=0$ without loss of generality. The second assumption is that $\phi$ is a increasing function on $[0,b]$, so that $\phi(0) < \phi(t)$ for any $t$ in the interval. With that we note that the coefficient $\ee^{-\lambda \phi(t)} / \ee^{-\lambda \phi(0)}$ decreases exponentially, that is, the integral is accumulated as the neighborhood of $0$, i.e.
\begin{equation}
	\Phi(\lambda) \approx \int_{0}^{\delta} g(t) \ee^{-\lambda \phi(t)} \dd t.
\end{equation}
As we know $\phi(t)$ to be increasing, its first non zero derivative must be of even order, so that when we expanded the function by Taylor,
\begin{equation}
\phi(t) \approx \phi(0) + \phi_{(2k+1)}(0) t^{2k+1},
\end{equation}
for some $k\geq 0$ and $|t| < \delta$ where $\phi_{(2k+1)} \deff \phi_{(2k+1)}/(2k+1)!$. As $\phi_{2k+1}$, the $(2k+1)$th-derivative, must be positive we also have that, as $\lambda \rightarrow \infty$,
	\begin{align}
		\Phi(\lambda) & \approx \ee^{-\lambda \phi(0)} \int_{0}^{\delta} g(t) \ee^{-\lambda \phi_{(2k+1)}(t) \cdot t^{2k+1}} \dd t,\\
		& \stackrel{(u=t(\lambda \phi_{(2k+1)}(0))^{1/(2k+1)})}{\approx} \ee^{-\lambda \phi(0)} \frac{g(0)}{(\lambda \phi_{(2k+1)}(0))^{\frac{1}{2k+1}}} \int_{0}^{\infty} \ee^{-u^{2k+1}} \dd u, \\
		& \stackrel{(s=u^{2k+1})}{\approx} \ee^{-\lambda \phi(0)} \frac{g(0)}{ (2k+1)(\lambda \phi_{(2k+1)}(0))^{\frac{1}{2k+1}}} \int_{0}^{\infty} \ee^{-s} \dd s.
	\end{align}
where the integral can be evaluated as a Gamma function. More than anything we need to focus on an important assumption that was made: that $\phi$ has a global minima at one of the endpoints of integration. If that were not the case, for example if $a < 0 < b$ and $\phi(t)$ had a global minima at $t=0$ we would now have to consider the expansion
\begin{equation}
\phi(t) \approx \phi(0) + \phi_{2k}(0) t^{2k}	
\end{equation}
which would lead us to the expression
\begin{equation}
\Phi(\lambda) \approx \ee^{-\lambda \phi(0)} \frac{g(0)}{(\lambda \phi_{2k}(0))^{\frac{1}{2k}}} \int_{-\infty}^{\infty} \ee^{-u^{2k}} \dd u.
\end{equation}
This differentiates of the previously result more importantly by the limits of integration. This is closely related to physical boundaries of the problem. Before enunciating the proper Laplace Method we state the Implicit Function Theorem, which will become important in a series of results that follows.

\begin{teorema}	\label{Theorem: Implicit Function}
	(Implicit Function Theorem) Assume that $S$ is an open subset of $\R^{2}$ and that $F: S \rightarrow \R$ is a function of class $C^1$. Assume also that $(a, b)$ is a point in $S$ such that
	\begin{equation}
		F(a, b) = 0, \quad \text{and} \quad D_yF(a,b)\neq 0.
	\end{equation}
	Then
	\begin{enumerate}
		\item There exist $r_0, r_1 > 0$ such hat for every $x \in \R$ such that $|x-a| < r_0$, there exists a unique $y \in \R$ such that $|y - b| < r_1$
		\begin{equation}
			F(x,y) = 0.
		\end{equation}
		That is, defines a function $y = f(x)$ for $x \in \R$ near $a$ and $y$ close to $b$.
		\item Moreover, the function $f: B(r_0, a) \rightarrow B(r_1, b) \subset \R$ from item $1$ is of class $C^1$, and its derivatives may be determined by differentiating the identity 
		\begin{equation}
			F(x, f(x)) = 0
		\end{equation}
		and solving to find the partial derivatives of $f$.
	\end{enumerate}
\end{teorema}

Now we turn to our main interest.

\begin{teorema}
	(Laplace's Method, endpoint contributions) Let $\Phi$ be as in \eqref{equation: Laplace}, with a finite and $g \in L^1(a,b)$. Assume that $\phi$ has a unique global minimum at $t=a$. In addition, suppose that for some $\delta > 0$ and some $\alpha > -1$, the function $g$ takes the form
	\begin{equation}
		g(t) = (t-a)^{\alpha} g_0(t)
	\end{equation}
	for $|t-a| < \delta$ and $g_0$ continuously differentiable in $[a, a+\delta]$ and $g_0(a) \neq 0$, whereas $\phi$ is twice continuously differentiable on $[a, a+\delta]$ with $\phi'(a) > 0$. Then the function $\Phi$ satisfies an estimate of the form
	\begin{equation}
		\Phi(\lambda) = \frac{1}{\lambda^{\alpha + 1}} \frac{g_0(a) \Gamma(\alpha + 1)}{\phi'(a)^{\alpha + 1}} \ee^{-\lambda \phi(a)} \left(1+\Boh(\lambda^{-1})\right)
	\end{equation}
	as $\lambda \rightarrow \infty$.
\end{teorema}

\begin{proof}
	Start by assuming, without loss of generality, that $a=0$ and dividing the integral in two intervals, one in the region where we have information about the behavior of $g$ and another in the rest of the interval,
	\begin{equation}
		\Phi(\lambda) = \int_{0}^{b} g(t) \ee^{-\lambda \phi(t)} \dd t = \int_{0}^{\delta} g(t) \ee^{-\lambda \phi(t)} \dd t + \int_{\delta}^{b} g(t) \ee^{-\lambda \phi(t)} \dd t = (1) + (2).
	\end{equation}
	Knowing that $\phi$ has a global minima at $t = 0$ we assume, changing $\delta$ as needed, that $\phi(t) > \phi(\delta)$ for any $t > \delta$. That give us a simple bound for $(2)$:
	\begin{equation}
		|(2)| \leq \int_{\delta}^{b} |g(t)| \ee^{-\lambda \phi(t)} \dd t \leq \ee^{-\lambda \phi(\delta)} \int_{\delta}^{b} |g(t)| \dd t \leq \ee^{-\lambda \phi(\delta)} \normL{1}{g}{[0,b]}.
	\end{equation}
	We now get back to $(1)$ and start by studying the difference given by $\phi(t) - \phi(0) = s$ in a similar way that we would think of a change of variable from $t$ to $s$. Define the multivariate function
	\begin{equation}
		H(t,s) \deff \phi(t) - \phi(0) - s
	\end{equation} 
	such that $H(0,0) = 0$ and $\partial_t H(0,0) \neq 0$. This means we are in the condition to use \autoref{Theorem: Implicit Function} to say that there exists a unique bijective and continuously differentiable function $t(s) \deff t$ such that 
	\begin{equation}
		\begin{cases}
			\phi(t(s)) - \phi(0) - s = 0,	\\
			t'(0) = \frac{\partial_s H(0,0)}{\partial_t H(0,0)} = \frac{1}{\phi'(0)}.
		\end{cases}
	\end{equation}
	If we set $T = t^{-1}(\delta)$ we would then have for $(1)$
	\begin{equation}
		\int_{0}^{\delta} g(t) \ee^{-\lambda \phi(t)} \dd t = \int_{0}^{T} g(t(s)) \ee^{-\lambda \phi(t(s))} t'(s) \dd s.
	\end{equation}
	Now, to apply the assumption on $g$ we write $g(t(s)) = t^\alpha(s) g_0(t(s))$. As $t$ is $C^2$ and vanishes at the origin, using its Taylor expansion, we have
	\begin{equation}
	t(s) = s t'(0) (1+r(s)) = \frac{s}{\phi'(0)}(1+r(s)).		
	\end{equation}
	Joining both results,
		\begin{align}
			\int_{0}^{T} g(t(s)) \ee^{-\lambda \phi(t(s))} t'(s) \dd s & = \ee^{\lambda \phi(0)} \int_{0}^{T} g(t(s)) \ee^{-\lambda (\phi(t(s)) - \phi(0))} t'(s) \dd s \\
			& = \ee^{\lambda \phi(0)} \int_{0}^{T} t^\alpha(s) g_0(t(s)) \ee^{-\lambda s} t'(s) \dd s \\
			& = \frac{\ee^{\lambda \phi(0)}}{\phi'(0)^\alpha} \int_{0}^{T} s^\alpha (1+r(s))^\alpha g_0(t(s)) \ee^{-\lambda s} t'(s) \dd s			
		\end{align}
	where, if we set $f(s) \deff s^\alpha (1+r(s))^\alpha g_0(t(s))$ we apply Watson's Lemma to finish the result. 
\end{proof}

\begin{teorema} 	\label{prop: (Laplace's Method, interior contributions)}
	(Laplace's Method, interior contributions) Let $\Phi$ be as in \eqref{equation: Laplace}, with a finite and $g \in L^1(a,b)$. Assume that $\phi$ has a unique global minimum at $t=c \in (a,b)$. In addition, suppose that for some $\delta > 0$ and some $\alpha > -1$, the function $g$ takes the form
	\begin{equation}
		g(t) = (t-a)^{\alpha} g_0(t)
	\end{equation}
	for $|t-a| < \delta$ and $g_0$ continuously differentiable in $[c-\delta, c+\delta]$ and $g_0(a) \neq 0$, whereas $\phi$ is three times continuously differentiable on $[c-\delta, c+\delta]$ with $\phi'(a) > 0$. Then the function $\Phi$ satisfies an estimate of the form
	\begin{equation}
		\Phi(\lambda) = \ee^{-\lambda \phi(c)} \left( \frac{\sqrt{2\pi} g(c)}{\sqrt{|\phi''(c)|}\lambda^{1/2}} + \Boh(\lambda^{-3/2})\right)
	\end{equation}
	as $\lambda \rightarrow \infty$.
\end{teorema}

\begin{proof}
	Although the proof may be similar there is an important distinction in the solution. Suppose we wanted to try the same approach. We would first assume, w.l.g., that c = 0. Then, we could write
	\begin{align}
	\Phi(\lambda) & = \int_{a}^{b} g(t) \ee^{-\lambda \phi(t)} \dd t \\ 
	& = \int_{a}^{-\delta} g(t) \ee^{-\lambda \phi(t)} \dd t + \int_{-\delta}^{\delta} g(t) \ee^{-\lambda \phi(t)} \dd t + \int_{\delta}^{b} g(t) \ee^{-\lambda \phi(t)} \dd t\\
	&  \revdeff (1) + (2) + (3).		
	\end{align}
	For $(1)$ and $(3)$ the result is as it was before
	\begin{equation}
		\begin{cases}
			|(1)| \leq \int_{\delta}^{a} |g(-t)| \ee^{-\lambda \phi(-t)} \dd t \leq \ee^{\lambda \phi(-\delta)} \normL{1}{g}{[a,b]},\\
			|(3)| \leq \int_{\delta}^{b} |g(t)| \ee^{-\lambda \phi(t)} \dd t \leq \ee^{\lambda \phi(\delta)} \normL{1}{g}{[a,b]}.
		\end{cases}
	\end{equation}
	For $(2)$ we would try to define a function
	\begin{equation}
		H(t,s) \deff \phi(t) - \phi(0) - s^2.		
	\end{equation} 
	However we could not apply \autoref{Theorem: Implicit Function} as the partial derivative $\partial_t H = \phi'$ vanishes at $t=0$. A nice fact about this indetermination is that it does not come from the fact that there are no solutions to $H(t,s) = 0$ but from the fact that there are two. To solve this we will use a technique called \textit{blow up}. First, we introduce a variable $v \deff v(t)$ such that $vs = t$. Now, we redefine our function of interest to be
	\begin{equation}
		L(s, v) \deff \frac{\phi(sv) - \phi(0)}{s^2} - 1.		
	\end{equation}
	for $s\in(-\delta, \delta) - \{0\}$. We would think this function is also singular but note that, using the Taylor Expansion of $\phi$ we write
	\begin{equation}
		\phi(t) = \frac{\phi''(0)}{2}t^2 + t^3R(t)
	\end{equation}	
	for some error function $R$, which leave us with 
	\begin{equation}
		L(s, v) = \frac{\phi''(0)}{2}v^2 + sv^3R(sv) - 1,
	\end{equation}
	and we can indeed apply the Implicit Function Theorem. In fact, we know $L \in C^2$ and as we want to solve for $L(s,v) = 0$, we need that the $v_0$ we take to satisfy $L(0, v_0) = 0$, ie,
	\begin{equation}
		v_0 = \sqrt{\frac{2}{\phi''(0)}}.
	\end{equation}
	Now, we also know that
	\begin{equation}
		\partial_v L(0, v_0) = \phi''(0) v_0 \neq 0.
	\end{equation}
	So we get a bijective and continuously differentiable function $v = v(s)$ such that
	\begin{equation}
		\begin{cases}
			L(s, v(s)) = 0 \\
			v(0) = v_0
		\end{cases}
	\end{equation}
	But, $L(s, v(s))$ means exactly that
	\begin{equation}
		\phi(t) - \phi(0) = s^2
	\end{equation}
	with $t = t(s) = sv(s)$. Getting back to our equation, we have that
	\begin{equation}
		\ee^{\lambda \phi(0)} \int_{-\delta}^{\delta} g(t) \ee^{-\lambda (\phi(t)-\phi(0))} \dd t = \ee^{\lambda \phi(0)} \int_{\alpha}^{\beta} g(sv(s)) (v(s) + sv'(s)) \ee^{-\lambda s^2} \dd s
	\end{equation}
	where we calculate $\alpha = - \sqrt{\phi(-\delta) - \phi(0)}$ and $\beta = \sqrt{\phi(\delta) - \phi(0)}$. One can now again finish the proof with \autoref{prop: cor watson}.
\end{proof}

\begin{exemplo} 	\label{Example: Hermite o(1)}
	\cite{nica2025gaussianintegralformulahermite}
	For $x \in \R$ fixed, i.e. $x = \Boh(1)$,  we have the following asymptotic as $n \rightarrow \infty$:
	\begin{equation}
		\He_n(x) = \sqrt{2} \left(\frac{n}{\ee}\right)^{n/2} \ee^{x^2/2} \cos{\left(\frac{n\pi}{2} - x \sqrt{n}\right)} (1 + \boh(1)).
	\end{equation}
	To see that, consider the integral representation of the Hermite Polynomial and re-write
	\begin{align}
		\He_n(x) & = \frac{1}{\sqrt{2\pi}} \int_{-\infty}^{\infty} (x + \ii t)^n \ee^{-t^2/2} \dd t = \frac{2}{\sqrt{2\pi}} \Re\left\{ \int_{0}^{\infty} (x + \ii t)^n \ee^{-t^2/2} \dd t \right\} \\
		& = \frac{2}{\sqrt{2\pi}} \Re\left\{ \int_{\Re(z) = x, \Im(z)>0} \ee^{n\ln{(z)} + \frac{1}{2}(z-x)^2} \frac{\dd z}{\ii} \right\} \\
		& = \frac{2}{\sqrt{2\pi}} \Re\left\{ \int_{\Re(z) = 0, \Im(z)>0} \ee^{n\ln{(z)} + \frac{1}{2}(z-x)^2} \frac{\dd z}{\ii} \right\}
	\end{align}
	where you start with a line integral in $\C$ and move the contour over to the positive imaginary axis - justified by the Cauchy's Theorem and the exponential decay of the integrand near infinity. Now, we focus on $n\ln{(z)} + \frac{1}{2}(z-x)^2$ which has maximum at $z^* = \ii \sqrt{n} + \boh(\sqrt{n})$. The idea is to change variables by a factor of $\sqrt{z}$ so that the location of the maximum contribution stays $\Boh(1)$ as $n$ gets large. Specifically, we do
	\begin{equation}
		z = \ii \sqrt{n} s, \quad \frac{\dd z}{\ii} = \sqrt{n} \dd s, \quad \ln{(z)} = \ln{(\sqrt{n})} + \ln{(s)} + \ii \frac{\pi}{2}
	\end{equation}  
	which give us
	\begin{equation}
		\frac{1}{2}(z-x)^2 = -\frac{1}{2} ns^2 + \frac{1}{2} x^2 - \ii \sqrt{n} sx	
	\end{equation}
	so that we remain with
	\begin{align}
		\He_n(x) & = \frac{2}{\sqrt{2\pi}} \Re\left\{ \int_{0}^{\infty} \ee^{n\ln{(\sqrt{n})}} \ee^{n\ln{(s)}} \ee^{n\ii\frac{\pi}{2}} \ee^{\frac{1}{2}x^2}\ee^{-\ii\sqrt{n}sx} \sqrt{n} \dd s \right\} \\
		& = \frac{2}{\sqrt{2\pi}} \sqrt{n} \ee^{n\ln{(\sqrt{n})}} \ee^{\frac{1}{2}x^2}  \int_{0}^{\infty} \ee^{n\left(\ln{(s)} - \frac{1}{2}s^2 \right)} \cos{\left(n\frac{\pi}{2} - \sqrt{n} s x\right)} \dd s,
	\end{align}
	where we can apply Laplace's approximation discussed before. Noting that $\phi(s) = \ln{(s)} - \frac{1}{2}s^2$ has $\phi'(s) = \frac{1}{s} - s$ and $\phi''(s) = -\frac{1}{s^2} - 1$, with global maximum at $s^* = 1$ with $\phi(1) = -\frac{1}{2}$ and $\phi''(1) = -2$, yielding
	\begin{align}
		\He_n(x) & = \frac{2}{\sqrt{2\pi}} \sqrt{n} \ee^{n\ln{(\sqrt{n})}} \ee^{\frac{1}{2}x^2}  \left( \ee^{n\phi(1)} \sqrt{\frac{2\pi}{n|\phi''(1)|}} \right) \cos{\left(n\frac{\pi}{2} - \sqrt{n}x\right)} (1+\boh(1)), \\
		& = \sqrt{2} \ee^{n \ln{(\sqrt{n})}} \ee^{\frac{1}{2}(x^2 - n)} \cos{\left(n\frac{\pi}{2} - \sqrt{n}x\right)} (1+\boh(1)), \\
		& \sqrt{2} \left(\frac{n}{\ee}\right)^{n/2} \ee^{x^2/2} \cos{\left(\frac{n\pi}{2} - x \sqrt{n}\right)} (1 + \boh(1)).
	\end{align}
\end{exemplo}

\subsection{Steepest Descent}

Let us now finally explore integrals over the complex plane with complex valued functions. We retake the integral 
\begin{equation}
	\phi(\lambda) := \int_{\Gamma} g(s) \ee^{-\lambda f(s)} \dd s,
	\label{eq: general int}
\end{equation}
where $\Gamma$ is a contour in the complex plane and $f, g$ are given complex-valued, analytic functions. Suppose that $\Gamma$ is parametrized by $\gamma : [a,b] \rightarrow \C$ and let us decompose the functions we have into their real and imaginary parts such as
\begin{equation}
	f(z) = u(z) + \ii v(z), \ \ g(z) = p(z) + \ii q(z), \quad \gamma(t) = \alpha(t) + \ii \beta(t)
\end{equation}
with all functions in the decomposition being real valued. Of course we now have 
\begin{equation}
	\ee^{-\lambda f(z)} = \ee^{-\lambda u(z)} \ee^{-\ii \lambda v(z)},
\end{equation}
and when $\lambda$ grows large, we have the exponentially growing/decaying term $\ee^{-\lambda u(z)}$ and a oscillating term coming from $\ee^{\ii \lambda v(z)}$ which is hard to control. The idea to get around this fact is to use paths in which this term is constant. If we have such a case, we would write
\begin{multline}
	\Phi(\lambda) = \ee^{-\ii \lambda v_0} \int_{a}^{b} (p(\gamma(t))\alpha'(t)-q(\gamma(t))\beta'(t)) \ee^{-\lambda u(\gamma(t))} \dd t \\
	+ \ii \ee^{-\ii \lambda v_0} \int_{a}^{b} (p(\gamma(t))\beta'(t)-q(\gamma(t))\alpha'(t)) \ee^{-\lambda u(\gamma(t))} \dd t
	\label{eq: decomposition}
\end{multline}
where each integral can now be solved by Laplace's method.

\begin{teorema} (Steepest Descent)
	Consider $\Phi$ as in \eqref{eq: general int} with a smooth contour $\Gamma$. Assume that the exponent $f = u + \ii v$ is an analytic function on a neighborhood of $\Gamma$ that satisfies
	\begin{equation}
		v(z) \equiv v_0, \quad z \in \Gamma
	\end{equation}
	In addition, suppose that $u$ has a unique minimum at an interior point of $\Gamma$, say on $z_0 \in \Gamma$, with $f''(z_0) \neq 0$, and also that $g$ is analytic on a neighborhood of $\Gamma$  and satisfies
	\begin{equation}
		\int_{\Gamma} |g(z)| |\dd z| < \infty.
	\end{equation}
	Then $\Phi$ has a leading asymptotic as $\lambda \rightarrow +\infty$ given by
	\begin{equation}
		\Phi(\lambda) = \ee^{-\lambda f(z_0)} \left(\frac{\sqrt{2\pi}g(z_0)\ee^{\ii \theta(z_0)}}{\lambda^{1/2}\sqrt{|f''(z_0)|}} + \Boh(\lambda^{-3/2}) \right), \ \ \lambda \rightarrow + \infty,
	\end{equation}
	where $\theta(z_0)$ is the angle of the oriented tangent of $\Gamma$ at $z_0$.
\end{teorema}

\begin{proof}
	Fix a parametrization $\gamma : [a,b] \rightarrow \C$ of $\Gamma$ with $\gamma(c) = z_0, a < c < b$, and decompose $\Phi$ as before in \eqref{eq: decomposition}. By the assumptions on the functions $f, g$ we apply Laplace`s method on both equations, and yields, as $\lambda \rightarrow + \infty$,
	\begin{equation}
		\int_{a}^{b} (p(\gamma(t))\alpha'(t)-q(\gamma(t))\beta'(t)) \ee^{-\lambda u(\gamma(t))} \dd t = \sqrt{2\pi} \ee^{-\lambda u(z_0)} \left( \frac{p(z_0)\alpha'(c)-q(z_0)\beta'(c)}{\sqrt{|(u \circ \gamma)''(c)|} \lambda^{1/2}} + \Boh \left(\lambda^{3/2}\right) \right)
	\end{equation}
	and
	\begin{equation}
		\int_{a}^{b} (p(\gamma(t))\beta'(t)-q(\gamma(t))\alpha'(t)) \ee^{-\lambda u(\gamma(t))} \dd t = \sqrt{2\pi} \ee^{-\lambda u(z_0)} \left( \frac{p(z_0)\beta'(c)-q(z_0)\alpha'(c)}{\sqrt{|(u \circ \gamma)''(c)|} \lambda^{1/2}} + \Boh \left(\lambda^{3/2}\right) \right)
	\end{equation}
	Using that $v_0 + u(z_0) = f(z_0)$ and also that,
	\begin{equation}
		 (p(z_0)\alpha'(c)-q(z_0)\beta'(c)) + \ii (p(z_0)\beta'(c)-q(z_0)\alpha'(c)) = g(z_0)\gamma'(c),
	\end{equation}
	and, combining the estimates for both parts we get
	\begin{equation}
		\Phi(\lambda) = \sqrt{2\pi} \ee^{-\lambda f(z_0)} \left( \frac{g(z_0)\gamma'(c)}{\sqrt{|(u \circ \gamma)''(c)|} \lambda^{1/2}} + \Boh \left(\lambda^{3/2}\right) \right)
	\end{equation} 
	and it suffices now to get an expression for the second derivative in the denominator. Knowing the imaginary part is constant we write
	\begin{equation}
		\frac{\dd}{\dd t} u(\gamma(t)) = \frac{\dd}{\dd t} f(\gamma(t)) = \gamma'(t) f'(\gamma(t))
	\end{equation}
	and
	\begin{equation}
		\frac{\dd^2}{\dd t} u(\gamma(t)) = \frac{\dd^2}{\dd t} f(\gamma(t)) = (\gamma'(t))^2 f''(\gamma(t)) + \gamma''(t) f'(\gamma(t)).
	\end{equation}
	Hence,
	\begin{equation}
		\frac{\gamma'(c)}{\sqrt{|(u \circ \gamma)''(c)|}} = \frac{1}{\sqrt{|f''(z_0)|}} \frac{\gamma'(c)}{|\gamma'(c)|}
	\end{equation}
	which concludes the proof.
\end{proof}

We now need do discuss why this assumption of constant value of $\Im(f(z))$ is not an assumption as strong as one might assume at first. Suppose we want to integrate in a arbitrary contour $\Gamma$ of the plane. Of course, our assumption that the functions were analytic permit us to deform the contour keeping the value of the integral. We are going to use this to find a contour $\tilde{\Gamma}$ such that sections of the contour are such that $\Im(f(z))$ is constant. We already know that being analytic means satisfying the \textit{Cauchy-Riemann} equations,
\begin{equation}
	u_x(z) = v_y(z), \quad \text{and} \quad u_y(z) = -v_x(z).
\end{equation}
An immediate result from these equations is that the gradients of $u$ and $v$ must be normal. Indeed, if $\langle \cdot, \cdot \rangle$ is the inner product on $\R^2$,
$$\langle \nabla u(z), \nabla v(z) \rangle = u_x(z)v_x(z) + u_y(z)v_y(z) = u_x(z)v_x(z) - v_x(z)u_x(z) = 0.$$
As we are following the curves where  $\Im(f(z))$ is constant, we are also following the path of Steepest Descent of the real part of the integral. As we also are looking to use Laplace's method on the integrals, we may want to have our paths going through the critical points of the function. This fully motivates the method described.


\begin{exemplo}
	There are three asymptotics regimes for $\He_n(x)$ when $x \rightarrow \infty$ depending on the value of $\lim_{N \rightarrow \infty} x/2\sqrt{N}$ for $N = n+1$, each of which has qualitatively different behavior
	\begin{itemize}
		\item Small $x$ values $-2\sqrt{N} < x < 2\sqrt{N}$: in this regime the Hermite Polynomials will be oscillating, and a have a similar qualitative behavior as in \autoref{Example: Hermite o(1)}, with modification to phase and amplitude.
		\item Large $x$ values $|x| > 2\sqrt{N}$: In this regime the polynomial does not have more oscillations and is monotone increasing. As $x$ grows larger, this approximation behaves as $x^N \ee^{-N^2/2x^2}$.
		\item Edge $x$ values $x = \pm 2\sqrt{N} + \Boh(N^{-\frac{1}{6}})$: in this regime the polynomials will be approximated by the Airy Function. This interpolates between the oscillatory small $x$ regime and the monotone large $x$ regimes. 
	\end{itemize} 
	
	We write the integral representation given by:
	\begin{equation}
		\begin{split}
			\He_n(x) & \deff \frac{n!}{2\pi \ii} \int_{\Gamma} z^{-n-1} \ee^{xz-\frac{z^2}{2}} \dd z \\
			& = \frac{n!}{2\pi \ii} \int_{\Gamma} \ee^{xz-\frac{z^2}{2} - (n+1)\ln{(z)}} \dd z
		\end{split}
	\end{equation}
	where $\Gamma$ is a contour that goes around the origin coming and going from minus infinity such as $\Gamma = \{\lambda - \ii \delta : \lambda < 0 \} \cup \{\delta \ee^{\ii \theta}: -\pi/2 < \theta < \pi/2 \} \cup \{\lambda + \ii \delta : \lambda < 0 \}$. This is equivalent to using $\Gamma = \{\ee^{\ii \theta}: 0 \leq \theta < 2\pi \}$.	To determine $\He_{n}$, set $N = n + 1$ and consider the scaling $x \mapsto N^a \alpha$ and $z = N^b \beta$ 
	\begin{equation}
		\begin{split}
			\He_n(N^a \alpha) & = \frac{(N-1)!}{2\pi \ii} \int_{\Gamma} \ee^{(N^a \alpha)(N^b \beta)- \frac{1}{2}(N^b \beta)^2 - N\ln{(N^b \beta)}} \dd \beta \\
			& \stackrel{(a = b = 1/2)}{=}   \frac{(N-1)! N^{\frac{1}{2}\left(1-N\right)}}{2\pi \ii} \int_{\Gamma}  \ee^{-N(- \alpha \beta+ \frac{1}{2} \beta^2 + \ln{(\beta)})} \dd \beta,
		\end{split}
	\end{equation}
	so that we have
	\begin{equation}
		\He_n(\sqrt{N}\alpha) = C_n^0 \int_{\Gamma}  \ee^{-N \phi_\alpha(\beta) } \dd \beta
		\label{Equation: Hermite Integral}
	\end{equation}
	with 
	\begin{equation}
		\phi_\alpha(\beta) = \frac{1}{2} \beta^2 - \alpha \beta + \ln{(\beta)}.
	\end{equation}
	%The same can be done in the second case, using now the re-parametrization, that is
	%\begin{equation}
	%	\begin{split}
		%		\He_{n-1}(\sqrt{N} \alpha) & \stackrel{(a = b = 1/2)}{=}   \frac{n! N^{-N/4}}{2\pi \ii} \int_{\Gamma}  \ee^{-N(-2 \alpha \beta+ \beta^2 + \frac{1}{2}\ln{(\beta)})} \beta^{-1} \dd \beta
		%	\end{split}
	%\end{equation}
	%So that we have
	%\begin{equation}	
	%	\He_{n-1}(\sqrt{N}\alpha) = C_n^1 \int_{\Gamma}  \ee^{-N \phi_\alpha(\beta)} \beta^{-1} \dd \beta.
	%\end{equation}
	%As noted, for both cases, we have the same defining function $\phi_\alpha(\beta)$.
	As we want to calculate the asymptotic by steepest descent, we need to determine the critical points of such a function and then the steepest descent paths passing through this points. We can easily determine the critical points as $\beta_c^\pm = \frac{\alpha \pm \sqrt{{\alpha}^2 - 4}}{2}$ - solutions to the quadratic formula given by solving the critical point of the function (the zeros of its derivatives). Now, the behavior of the critical points will depend on the interval of the real parameter $\alpha$,
	\begin{equation}
		\beta_c^{\pm} = \frac{\alpha \pm \sqrt{\alpha^2 - 4}}{2} =
		\begin{cases}
			\begin{aligned}
				\ee^{\pm \gamma}, \quad & \alpha = 2 \cosh{(\gamma)} > 2, \quad \gamma > 0; \\
				\ee^{\pm \ii \gamma}, \quad & \alpha = 2 \cos{(\gamma)} \in [0,2), \quad \gamma \in (0,\pi/2]; \\
				1, \quad & \alpha = 2,
			\end{aligned}
		\end{cases}
	\end{equation}
	each of the cases refers to one of the regimes we described before.
		\begin{figure}[H]
		\begin{subfigure}{0.3\textwidth}
		\centering
		\begin{tikzpicture}[line cap=round,line join=round,>=triangle 45,x=1cm,y=1cm]
			\begin{axis}[
				title={$\alpha > 2$},
				xmin=-2, xmax=2,
				ymin=-2, ymax=2,
				xlabel={$x$},
				ylabel={$y$},
				trig format plots=rad,
				axis equal]					
				\node at (axis cs:0.7,0.7) {$t_{c_+}$};
				\node at (axis cs:-0.7,0.7) {$t_{c_-}$};
				\node at (axis cs:0.433, 0.25) {\textbullet};
				\node at (axis cs:-0.433,0.25) {\textbullet};
				\addplot[line width=1pt,samples=100,domain=pi/3:pi] 
				(
				{
					((cos(3.141592653589793/3)-sqrt((cos(3.141592653589793/3))^(2)-(x-3.141592653589793/3+sin(2*3.141592653589793/3)/2)/tan(x)))/(2*cos(x)))*sin(x)
				},
				{
					((cos((3.141592653589793/3))-sqrt((cos(3.141592653589793/3))^(2)-(x-3.141592653589793/3+sin(2*3.141592653589793/3)/2)/tan(x)))/(2*cos(x)))*cos(x)
				}
				);
				\addplot[line width=1pt,samples=100,domain=0:pi/3] 
				(
				{
					((cos(3.141592653589793/3)+sqrt((cos(3.141592653589793/3))^(2)-(x-3.141592653589793/3+sin(2*3.141592653589793/3)/2)/tan(x)))/(2*cos(x)))*sin(x)
				},
				{
					((cos((3.141592653589793/3))+sqrt((cos(3.141592653589793/3))^(2)-(x-3.141592653589793/3+sin(2*3.141592653589793/3)/2)/tan(x)))/(2*cos(x)))*cos(x)
				}
				);
				\addplot[line width=1pt,samples=100,domain=pi/3:pi/2] 
				(
				{
					((cos(3.141592653589793/3)+sqrt((cos(3.141592653589793/3))^(2)-(x-3.141592653589793/3+sin(2*3.141592653589793/3)/2)/tan(x)))/(2*cos(x)))*sin(x)
				},
				{
					((cos((3.141592653589793/3))+sqrt((cos(3.141592653589793/3))^(2)-(x-3.141592653589793/3+sin(2*3.141592653589793/3)/2)/tan(x)))/(2*cos(x)))*cos(x)
				}
				);
				\addplot[line width=1pt,samples=100,domain=.614:pi/3] 
				(
				{
					((cos(3.141592653589793/3)-sqrt((cos(3.141592653589793/3))^(2)-(x-3.141592653589793/3+sin(2*3.141592653589793/3)/2)/tan(x)))/(2*cos(x)))*sin(x)
				},
				{
					((cos((3.141592653589793/3))-sqrt((cos(3.141592653589793/3))^(2)-(x-3.141592653589793/3+sin(2*3.141592653589793/3)/2)/tan(x)))/(2*cos(x)))*cos(x)
				}
				);
				\addplot[line width=1pt,samples=100,domain=pi/3:pi] 
				(
				{
					-((cos(3.141592653589793/3)-sqrt((cos(3.141592653589793/3))^(2)-(x-3.141592653589793/3+sin(2*3.141592653589793/3)/2)/tan(x)))/(2*cos(x)))*sin(x)
				},
				{
					((cos((3.141592653589793/3))-sqrt((cos(3.141592653589793/3))^(2)-(x-3.141592653589793/3+sin(2*3.141592653589793/3)/2)/tan(x)))/(2*cos(x)))*cos(x)
				}
				);
				\addplot[line width=1pt,samples=100,domain=0:pi/3] 
				(
				{
					-((cos(3.141592653589793/3)+sqrt((cos(3.141592653589793/3))^(2)-(x-3.141592653589793/3+sin(2*3.141592653589793/3)/2)/tan(x)))/(2*cos(x)))*sin(x)
				},
				{
					((cos((3.141592653589793/3))+sqrt((cos(3.141592653589793/3))^(2)-(x-3.141592653589793/3+sin(2*3.141592653589793/3)/2)/tan(x)))/(2*cos(x)))*cos(x)
				}
				);
				\addplot[line width=1pt,samples=100,domain=pi/3:pi/2] 
				(
				{
					-((cos(3.141592653589793/3)+sqrt((cos(3.141592653589793/3))^(2)-(x-3.141592653589793/3+sin(2*3.141592653589793/3)/2)/tan(x)))/(2*cos(x)))*sin(x)
				},
				{
					((cos((3.141592653589793/3))+sqrt((cos(3.141592653589793/3))^(2)-(x-3.141592653589793/3+sin(2*3.141592653589793/3)/2)/tan(x)))/(2*cos(x)))*cos(x)
				}
				);
				\addplot[line width=1pt,samples=100,domain=.614:pi/3] 
				(
				{
					-((cos(3.141592653589793/3)-sqrt((cos(3.141592653589793/3))^(2)-(x-3.141592653589793/3+sin(2*3.141592653589793/3)/2)/tan(x)))/(2*cos(x)))*sin(x)
				},
				{
					((cos((3.141592653589793/3))-sqrt((cos(3.141592653589793/3))^(2)-(x-3.141592653589793/3+sin(2*3.141592653589793/3)/2)/tan(x)))/(2*cos(x)))*cos(x)
				}
				);
			\end{axis}
		\end{tikzpicture}
		\end{subfigure} \hspace{1cm}
		\begin{subfigure}{0.3\textwidth}
			\centering
			\begin{tikzpicture}[line cap=round,line join=round,>=triangle 45,x=1cm,y=1cm]
				\begin{axis}[
					title={$\alpha \in [0,2)$},
					xmin=-2, xmax=2,
					ymin=-2, ymax=2,
					xlabel={$x$},
					yticklabel=\empty,
					trig format plots=rad,
					axis equal]
					\node at (axis cs:0,1.0) {$t_{c_+}$};
					\node at (axis cs:0,1/2-0.1) {$t_{c_-}$};
					\node at (axis cs:0,1.32) {\textbullet};
					\node at (axis cs:0,0.14) {\textbullet};
					\addplot[line width=1pt,samples=100,domain=0.01:pi] 
					(
					{
						(x/(2*sin(x)))/(cosh(1) + sqrt(cosh(1)^2 - x/tan(x))) * sin(x)
					},
					{
						(x/(2*sin(x)))/(cosh(1) + sqrt(cosh(1)^2 - x/tan(x))) * cos(x)
					}
					);
					\addplot[line width=1pt,samples=100,domain=0.001:pi/3] 
					(
					{
						(cosh(1) + sqrt(cosh(1)^2 - x/tan(x)))/(2*cos(x)) * sin(x)
					},
					{
						(cosh(1) + sqrt(cosh(1)^2 - x/tan(x)))/(2*cos(x)) * cos(x)
					}
					);
					\addplot[line width=1pt,samples=100,domain=0.01:pi] 
					(
					{
						-(x/(2*sin(x)))/(cosh(1) + sqrt(cosh(1)^2 - x/tan(x))) * sin(x)
					},
					{
						(x/(2*sin(x)))/(cosh(1) + sqrt(cosh(1)^2 - x/tan(x))) * cos(x)
					}
					);
					\addplot[line width=1pt,samples=100,domain=0.001:pi/3] 
					(
					{
						-(cosh(1) + sqrt(cosh(1)^2 - x/tan(x)))/(2*cos(x)) * sin(x)
					},
					{
						(cosh(1) + sqrt(cosh(1)^2 - x/tan(x)))/(2*cos(x)) * cos(x)
					}
					);
				\end{axis}
			\end{tikzpicture}
		\end{subfigure} \hfill
		\begin{subfigure}{0.3\textwidth}
			\centering
			\begin{tikzpicture}[line cap=round,line join=round,>=triangle 45,x=1cm,y=1cm]
				\begin{axis}[
					title={$\alpha = 2$},
					xmin=-2, xmax=2,
					ymin=-2, ymax=2,
					xlabel={$x$},
					yticklabel=\empty,
					trig format plots=rad,
					axis equal]
					\addplot[line width=1pt,samples=100,domain=0.01:pi-0.01] 
					(
					{
						(x/(2*sin(x)))/(1 + sqrt(1 - x*cot(x))) * sin(x)
					},
					{
						(x/(2*sin(x)))/(1 + sqrt(1 - x*cot(x))) *cos(x)
					}
					);
					\addplot[line width=1pt,samples=100,domain=-pi/2:-0.01] 
					(
					{
						(1 + sqrt(1 - x*cot(x)))/(2*cos(x)) * sin(x)
					},
					{
						(1 + sqrt(1 - x*cot(x)))/(2*cos(x)) * cos(x)
					}
					);
					\addplot[line width=1pt,samples=100,domain=-pi+0.01:-0.01] 
					(
					{
						(x/(2*sin(x)))/(1 + sqrt(1 - x*cot(x))) * sin(x)
					},
					{
						(x/(2*sin(x)))/(1 + sqrt(1 - x*cot(x))) *cos(x)
					}
					);
					\addplot[line width=1pt,samples=100,domain=0.01:pi/2-0.01] 
					(
					{
						(1 + sqrt(1 - x*cot(x)))/(2*cos(x)) * sin(x)
					},
					{
						(1 + sqrt(1 - x*cot(x)))/(2*cos(x)) * cos(x)
					}
					);
					\node at (axis cs:0,1) {$t_{c_{\pm}}$};
					\node at (axis cs:0,1/2) {\textbullet};
				\end{axis}
			\end{tikzpicture}
		\end{subfigure} \hfill
		\caption{The lines represent the solution to $\Im[f(z) - f(t_c)] = 0$ for the three possible cases of the Hermite Polynomial. These are the curves we use to perform the asymptotic analysis on the integral. Note the case $\alpha = 2$; the coalescence of critical point gives origin to the Airy function.}
	\end{figure}
\end{exemplo}